\part{Securing Spaces}

Unit 1 dealt with attacks that arrive over a screen. Unit 2 deals with attacks that arrive through
a door, and it starts by giving the rest of the book its central vocabulary:
Chapter~\ref{ch:2.1} introduces the CIA triad -- confidentiality, integrity, and availability --
along with the language of threat actors and risk that every later chapter, physical or digital,
will reuse. It is worth slowing down on this chapter even though it feels abstract, because ``is
this an integrity problem or a confidentiality problem'' turns out to be the first useful question
to ask about almost any incident in this course.

The remaining three chapters apply that vocabulary to a building. Chapter~\ref{ch:2.2} walks
through how physical security actually fails in practice -- tailgating through a badge-locked
door, an unattended workstation, a dumpster full of paper that should have been shredded --
because most physical breaches do not look like a movie heist. Chapter~\ref{ch:2.3} covers the
controls organizations put in place to stop those failures: badges, locks, mantraps, visitor logs,
and the access-control philosophy of least privilege that will reappear in Unit 5 applied to
software instead of hallways. Chapter~\ref{ch:2.4} closes the unit with detection -- cameras,
sensors, and alarms -- on the premise that prevention eventually fails and an organization needs
to know when it did.

A theme worth watching across this unit: nearly every physical control has a direct digital
analogue, and nearly every digital control has a physical one. Keep that pairing in mind, since
it is the same pattern this book will return to in every unit that follows.
